\subsection{The graded torsion-zero estimate used by the moduli construction}
\label{subsec:ftb-torsion-zero-estimate}

We give the graded argument used by the source, including the
finiteness and projective-point conditions needed in its application.
All fields in this subsection may have positive characteristic.

For a graded \(K\)-algebra \(T=\bigoplus_{e\ge0}T_e\), finite
as a module over \(K[x_1,\ldots,x_m]\), where \(x_i\in T_1\),
and a homogeneous ideal \(I\), write \(P_I\) for the eventual
polynomial of
\[
                  e\longmapsto\dim_K(T/I)_e.
\]
When \(\deg P_I\le r\), set
\[
                         D_r(I)=r![T^r]P_I(T).
 \tag{TZ1}\label{eq:ftb-graded-degree}
\]
The polynomial and the grading here are those of the actual finite
graded algebra, not an independently chosen projective embedding.

\begin{lemma}[The graded degree calculation]
\label{lem:ftb-graded-degree-calculation}
Suppose \(K\) is a field, \(T\) is as above, and
\(Q\) is a homogeneous prime with \(\dim(T/Q)=r+1\).
Then \(P_Q\) has degree \(r\) and \(D_r(Q)\) is a positive
integer.
For a homogeneous \(f\notin Q\) of degree \(a\), the
minimal primes \(Q'\) of \(Q+(f)\) are homogeneous and have
dimension \(r\); for \(r\ge1\),
\[
             \sum_{Q'\in\operatorname{Min}(Q+(f))}
                        D_{r-1}(Q')\le aD_r(Q).
 \tag{TZ2}\label{eq:ftb-graded-cut-degree}
\]
Finally, if \(\dim_KT_e\le(Ce)^g\) for every \(e\ge1\), then
\[
                             D_g(0)\le g!C^g.
 \tag{TZ3}\label{eq:ftb-graded-growth-degree}
\]
\end{lemma}
\begin{proof}
The finite module structure over \(K[X_1,\ldots,X_m]\),
where \(X_i\) acts by \(x_i\), gives finite-dimensional
graded pieces and their eventual Hilbert polynomial.
One can obtain this in exactly the form used here by extending
the grading by zero in negative degrees and applying
successive multiplication-by-\(X_i\) exact sequences to
a finite homogeneous module presentation.
For a homogeneous \(I\) and homogeneous \(f\) of degree \(a\),
the kernel and cokernel of multiplication by \(f\) give
\[
 P_I(T)=P_{I+(f)}(T)+P_{(I:f)}(T-a).
 \tag{TZ4}\label{eq:ftb-colon-hilbert-identity}
\]
Indeed, in degree \(e+a\) this is the dimension formula
for multiplication \( (T/(I:f))_e\to(T/I)_{e+a}\).
For prime \(Q\) with \(f\notin Q\), it becomes
\[
                   P_{Q+(f)}(T)=P_Q(T)-P_Q(T-a).
 \tag{TZ5}\label{eq:ftb-prime-hilbert-difference}
\]

Here are the dimension and coefficient arguments in the source.
Every quotient is finite over a quotient of the polynomial ring,
so its dimension is finite.
If all \(x_i\) lie in \(I\), then \(P_I=0\): its finite
graded module presentation is then annihilated by every positive
polynomial variable, and has only finitely many nonzero degrees.
For homogeneous primes avoiding one \(x_i\), the polynomial
\(P_Q\) is nonzero, since no power of that \(x_i\) lies in \(Q\).
Noetherian induction, using \eqref{eq:ftb-colon-hilbert-identity},
shows that \(\dim(T/I)\le s\) implies \(\deg P_I<s\).
For the prime induction step choose a degree-one \(x_i\notin Q\)
and cut by it; for a nonprime ideal choose a homogeneous
zero divisor witnessing nonprimality, so that both \(I+(f)\)
and \((I:f)\) strictly contain \(I\).
The dimension of every minimal component of a nonzero principal
cut of an affine domain is one less than that of the domain;
this is the affine finite-type dimension theorem and principal
ideal theorem, applied to \(T/Q\).
Using such a component and the inequality
\(\deg P_J\le\deg P_I\) for \(I\subset J\), the reverse
induction proves \(\deg P_Q=r\).
The dimension inequality just mentioned follows directly by
comparing the eventual dimensions of the quotient pieces.

For completeness, the coefficient comparison does not discard
multiplicities by assuming reducedness of a cut.
For any homogeneous \(I\) with \(\deg P_I\le s\),
Noetherian induction gives
\[
       \sum_{Q\in\operatorname{Min}(I)}[T^s]P_Q
                                     \le[T^s]P_I.
 \tag{TZ6}\label{eq:ftb-minimal-prime-coefficient}
\]
For prime \(I\) this is immediate.
Otherwise choose homogeneous \(f\) with both overideals
in \eqref{eq:ftb-colon-hilbert-identity} strictly larger.
Separate the minimal primes of \(I\) according as they contain
\(f\); those containing \(f\) are minimal over \(I+(f)\),
and the others are minimal over \((I:f)\).
Apply induction to the two collections and use
\eqref{eq:ftb-colon-hilbert-identity}; translation \(T\mapsto T-a\)
does not alter the coefficient of degree \(s\).
This proves \eqref{eq:ftb-minimal-prime-coefficient}.
The leading coefficient is nonnegative because the polynomial
is eventually the dimension of a vector space.
Moreover \(s![T^s]P_I\) is an integer whenever \(\deg P_I\le s\):
the \(s\)-fold finite difference of the eventual
integer-valued polynomial has that constant value.
For a prime of dimension \(s+1\) it is nonzero, hence at least one.
Equation \eqref{eq:ftb-prime-hilbert-difference}, followed by
\eqref{eq:ftb-minimal-prime-coefficient}, now gives
\eqref{eq:ftb-graded-cut-degree}.
Comparing the eventual polynomial \(P_0(e)\) with \(C^ge^g\)
gives \eqref{eq:ftb-graded-growth-degree} by leading coefficients.
\end{proof}

A graded character below means a \(K\)-algebra homomorphism
\(\phi:T\to K\).  Two such characters are \emph{projectively
equivalent} if there is \(c\in K\) such that
\[
                \phi'(y)=c^e\phi(y)
                         \quad(y\in T_e,\ e\ge0).
 \tag{TZ7}\label{eq:ftb-graded-scaling}
\]
We only count characters on which some \(x_i\) is nonzero;
then an equivalence between two counted characters has \(c\ne0\).

\begin{lemma}[The actual graded B\'ezout bound]
\label{lem:ftb-actual-graded-bezout}
Let \(K\) be algebraically closed and \(T\) a graded domain
finite over \(K[x_1,\ldots,x_m]\), \(x_i\in T_1\).
Suppose \(\dim T=g+1\), \(g\ge1\), and
\(\dim_KT_e\le(Ce)^g\) for \(e\ge1\).
Let \(0\ne s\in T_d\) and let \(\mathcal F\subset T_a\).
Assume that characters annihilating \(\mathcal F\) and nonzero
on some \(x_i\) have only finitely many equivalence classes
under \eqref{eq:ftb-graded-scaling}.
Then the number of distinct such classes annihilating \(s\) is at most
\[
                              C^g g!d\,a^{g-1}.
 \tag{TZ8}\label{eq:ftb-actual-graded-bezout-bound}
\]
\end{lemma}
\begin{proof}
Choose finitely many representative characters \(\Phi\).
If a homogeneous prime \(Q\) of dimension at least two
contained all of \(\mathcal F\), choose
\[
                              h=\sum_i t^ix_i\in T_1
\]
with \(\psi(h)\ne0\) for every \(\psi\in\Phi\)
nonzero on some \(x_i\).
Such \(t\) exists: each corresponding polynomial
\(\sum_i\psi(x_i)T^i\) is nonzero, and \(K\) is infinite.
The ideal \(Q+(h)\) is proper, by taking the degree-zero
part if it contained \(1\).
A minimal prime \(Q'\) over it has positive dimension:
if \(h\in Q\), take \(Q'=Q\); otherwise use the
principal-cut dimension equality.
Not all \(x_i\) lie in \(Q'\), since finiteness over
\(K[x_i]\) would otherwise make \(T/Q'\) zero-dimensional.
The affine Nullstellensatz supplies a character \(\phi\)
annihilating \(Q'\) but nonzero on one \(x_i\).
It annihilates \(\mathcal F\), so is equivalent to some
\(\psi\in\Phi\), with nonzero scaling factor.
Then \(\phi(h)\ne0\), contradicting \(h\in Q'\).
Consequently, for every such \(Q\), some \(f\in\mathcal F\)
lies outside \(Q\).

We next count homogeneous primes containing \(\mathcal F\)
and avoiding at least one \(x_i\).
Inside a homogeneous prime \(Q\) of dimension \(k+1\),
the number in any finite collection of them is at most
\[
                               a^kD_k(Q).
 \tag{TZ9}\label{eq:ftb-count-inside-prime}
\]
For \(k=0\) each is \(Q\).
Indeed, a strict containment gives a homogeneous element
outside \(Q\); a minimal component of its cut has dimension
zero and lies under the larger prime, whereas a homogeneous
prime avoiding some \(x_i\) has positive dimension.
Thus the collection has size at most one, bounded by \(D_0(Q)\).
For \(k\ge1\), choose \(f\in\mathcal F\setminus Q\) as above.
Each prime in the collection contains a minimal prime of
\(Q+(f)\).
Apply induction to each such component and sum.
Equation \eqref{eq:ftb-graded-cut-degree} gives
\eqref{eq:ftb-count-inside-prime}; overlap only decreases
the required cardinality.

For a character \(\phi\), put
\(P_\phi=\bigoplus_e(\ker\phi\cap T_e)\).
This homogeneous core is prime: the lowest nonzero graded
terms of two elements in the graded quotient show that
their product cannot vanish.
Distinct counted projective classes give distinct cores.
Indeed, if \(P_\phi=P_{\phi'}\) and \(\phi(x_i)\ne0\),
the homogeneous relation
\[
                    \phi(x_i)^e y-\phi(y)x_i^e\in P_\phi
\]
gives
\(\phi'(y)=(\phi'(x_i)/\phi(x_i))^e\phi(y)\).
Every core being counted contains \(s\), hence contains
a minimal prime \(Q'\) of \((s)\).
Apply \eqref{eq:ftb-count-inside-prime} with \(k=g-1\)
to these components, then \eqref{eq:ftb-graded-cut-degree}
to the cut of \(0\) by \(s\), and finally
\eqref{eq:ftb-graded-growth-degree}.
The result is
\(a^{g-1}dD_g(0)\le C^gg!d a^{g-1}\).
This proves the bound for every finite collection, hence
for the full finite set of classes.
\end{proof}

\paragraph{Why the finite graded algebra exists in the application.}
Let \(X/K\) be integral and \(M\) a line bundle.
The source constructs its section ring by writing a section of
\(M^e\) as compatible functions in its actual trivializing charts,
then placing that family as the coefficient of \(Z^e\)
in a polynomial ring.
The transition functions multiply under addition of exponents,
so these monomials form a graded algebra.
Restriction to any nonempty chart is injective on each
homogeneous piece: a section zero on that chart is zero on
all charts by integrality.
Mapping its coefficients to one such chart therefore embeds
the entire graded algebra in a domain of polynomials.
This proves that the section ring used above is a domain
with the actual section spaces as its graded pieces.

Suppose \(x_i\in H^0(X,M)\) is a very ample system.
On the affine chart \(x_i\ne0\), its ratios \(x_j/x_i\)
generate the coordinate ring.
For a homogeneous section \(y\) of degree \(e\), express
\(y/x_i^e\) as a polynomial in these ratios and homogenize:
\[
                          yx_i^n=G(x),\qquad\deg G=n+e.
 \tag{TZ10}\label{eq:ftb-section-ring-clearing}
\]
In the actual counting application take the subalgebra
\(T\) generated by the finitely many \(x_i\), the section
\(s\), and the finitely many equations of the relevant
finite fibre.
Give it their section degrees.
It is finite over \(S=K[x_i]\), as the following source
argument shows; this step is not inferred just from finite
generation as a \(K\)-algebra.

For homogeneous \(y\in T_e\), \eqref{eq:ftb-section-ring-clearing}
gives \(yx_i^{n_i}\in S\) for every \(i\).
The ideal
\[
                       \{b\in S:yb\in S\}
\]
contains a power of each \(x_i\), so contains a power
\((x_i)^N\) of their finitely generated ideal.
Let \(V\) be the \(S\)-span of the products of \(N\)
of the \(x_i\).
It is finite, nonzero, and faithful because \(T\) is a
domain and some \(x_i\ne0\).
The homogeneous component of \(S\) of degree \(e+N\)
is spanned by products of \(e+N\) of the \(x_i\).
Splitting each product after its first \(e\) factors
shows \(yV\subset V\).
The determinant trick gives a monic equation for \(y\)
over \(S\).
Sums of homogeneous integral elements are integral,
and \(T\) is of finite type over \(S\); hence \(T\) is
finite over \(S\).

Its dimension is \(g+1\) when the affine chart of \(X\)
has dimension \(g\).
Explicitly, let \(C\) be that chart ring, put \(r_i=1\)
and \(r_j=x_j/x_i\), and map
\(K[X_j]\to C[Z]\) by \(X_j\mapsto r_jZ\).
Homogenization shows its kernel equals the kernel of
\(X_j\mapsto x_j\).
Its image contains \(Z\), and every element of \(C[Z]\)
enters the image after multiplication by a power of \(Z\).
Thus the image and \(C[Z]\) have the same fraction field.
Finiteness of \(T/S\) now gives
\(\dim T=\operatorname{trdeg}_K C+1=g+1\).
The dimensions of the pieces of \(T\) are bounded by
those of the actual section spaces.

Finally a character of \(T\) nonzero on \(x_i\) gives
normalized coordinates \(\phi(x_j)/\phi(x_i)\).
All homogeneous relations hold there.
Homogenizing arbitrary chart relations yields a
\(K\)-algebra map from the affine chart ring to \(K\),
and hence an actual point of \(X(K)\).
Equation \eqref{eq:ftb-section-ring-clearing} implies
that its evaluation character is equivalent to \(\phi\)
on every degree.
Very ampleness also proves that distinct actual points
give distinct classes.
Therefore a finite geometric fibre really does provide
the finite set \(\Phi\) required in the preceding lemma.

\begin{proposition}[Zeros on prime-to-characteristic torsion]
\label{prop:ftb-torsion-zero-growth}
Let \(A/K\) be an abelian variety of dimension \(g\ge1\)
over an algebraically closed field.
Let \(M\) be a symmetric line bundle with a specified
very ample system, and suppose
\[
                h^0(A,M^e)\le(Ce)^g\qquad(e\ge1).
\]
For \(k\ge1\) prime to \(\operatorname{char}K\) and
\(0\ne s\in H^0(A,M^d)\),
\[
             \#\{x\in A[k](K):s(x)=0\}
                         \le C^g g!d\,k^{2g-2}.
 \tag{TZ11}\label{eq:ftb-torsion-zero-bound}
\]
\end{proposition}
\begin{proof}
Choose an index \(i_0\) with \(x_{i_0}(0)\ne0\).
The sections
\[
                    x_j-\frac{x_j(0)}{x_{i_0}(0)}x_{i_0}
\]
have common zero locus precisely the origin, by the affine
ratios of the very ample system.
Pull them back by \([k]\).
Symmetry gives the actual line-bundle isomorphism
\([k]^*M\simeq M^{k^2}\), so these pullbacks are equations
of degree \(a=k^2\) with common zero set \(A[k](K)\).
Their common zero set is finite, of cardinality \(k^{2g}\),
since \([k]\) is etale of degree \(k^{2g}\).
Use the finite graded subalgebra and point comparison
just constructed and apply
Lemma~\ref{lem:ftb-actual-graded-bezout}.
This gives \(C^gg!d(k^2)^{g-1}\), as claimed.
No torsion point outside this same marked fibre or
different line bundle enters the calculation.
\end{proof}

\paragraph{Source concordance.}
The section-ring construction is
\texttt{ThetaFamilyOrigin/ExistsSectionRing.lean}
(\texttt{ccbcde80e4e6666c2403ef1fd456ca5377bd4ab3}).
The clearing, finite-module and dimension arguments are
\texttt{ExistsMulMemAdjoinOfMemPowRange.lean}
(\texttt{fb2e32e96e5edc160ce0474fc8a810e096a2b3c3}),
\texttt{NonemptyGradedAlgebraSubGrading.lean}
(\texttt{b5b4b0c004dbe53fbf7722321c51fe5ac5f3bb5c}),
and \texttt{CardZerosOnTorsionCosetLe.lean}
(\texttt{72d5e39854bcd931cf08432e3f914729fde60834}).
The Hilbert-degree induction and finite-core counting are
\texttt{RingKrullDimQuotientOfMemMinimalPrimesSup.lean}
(\texttt{d1f4d4c443e5ba59b1d7022efe2848460942bc94})
and \texttt{RingKrullDimQuotientLeOne.lean}
(\texttt{80400608125d31abc9a5a9cdf23e09cd85097b5e}).
The final supplied-system application is
\texttt{ModuliProofs/CardZerosOnTorsionLeOfSymmetricVeryAmple.lean}
(\texttt{221898a2feed4bb454a82acc68c4e320ac130be9}).
These are all under
\texttt{src/Development/AbelianVarietyTheory/Proofs/}.
The remaining fine-moduli joins are the actual theta
very-ampleness/periodicity, the ideal Hilbert polynomial,
and the marked graph systems; this proposition does
not presume them.
